\begin{pastexam}[label={px:ntu-107-dsa-03}]{臺大資工所 107 學年度 第 III 題}
（單選；答錯倒扣當題分數）

(8 points) There are various collision resolution strategies for hash tables. Please answer the
following questions under the assumption of simple uniform hashing and all keys are unique, given
a hash table $H$ with $m$ slots to store $n$ elements.
\begin{enumerate}[label=\textbf{\arabic*.}, start=11]
\item (3 points) Suppose we use chaining with the hash function $h(k)=k\ \mathit{mod}\ 5$, and
insert the following 13 keys into $H$ with 5 slots in the exact sequence:
$32,39,21,3,6,40,34,24,15,25,7,13,50$. What is the expected number of key comparisons when
searching for an element with the key $k$? (A) 1.0 (B) 2.6 (C) 3.0 (D) 3.6.
\item (5 points) Suppose we use double hashing of the form $h(k,i)=(h_1(k)+ih_2(k))\mathit{mod}\ m$,
where the primary hash function is defined as $h_1(k)=k\ \mathit{mod}\ 7$ and the secondary hash
function is defined as $h_2(k)=5-(k\ \mathit{mod}\ 5)$. Given $H$ with 7 slots to insert the
following 4 keys in the exact sequence: $35,42,3,21$. What is the index of the slot storing the
element with the key 21? (A) 0 (B) 3 (C) 4 (D) 5.
\end{enumerate}
\end{pastexam}

\begin{pastsoln}
答案：11.~(B) 2.6\quad 12.~(C) 4。

\textbf{11.} 題目沒有說 $k$ 在不在表裡。 先把表建出來（新的鍵插在串列最前面）：
\begin{center}
\small
\begin{tabular}{cll}
\toprule
格 & 落在這一格的鍵（插入順序） & 串列（由頭到尾） \\
\midrule
0 & $40, 15, 25, 50$ & $50 \to 25 \to 15 \to 40$ \\
1 & $21, 6$ & $6 \to 21$ \\
2 & $32, 7$ & $7 \to 32$ \\
3 & $3, 13$ & $13 \to 3$ \\
4 & $39, 34, 24$ & $24 \to 34 \to 39$ \\
\bottomrule
\end{tabular}
\end{center}
負載因子 $\alpha = n/m = 13/5 = 2.6$。
\begin{itemize}
\item \textbf{$k$ 不在表裡}（失敗的搜尋）：在簡單均勻雜湊下，$k$ 落在每一格的機率都是 $1/5$，要把
那一條串列整條比完。 期望比較次數是串列的平均長度 $(4 + 2 + 2 + 2 + 3)/5 = 13/5 = 2.6$，也就是
$\alpha$（資結書定理~12.2 第 2 點）。
\item \textbf{$k$ 在表裡}（成功的搜尋）：每個鍵被找的機會相同時，長度 $L$ 的串列裡的鍵分別要比
$1, 2, \ldots, L$ 次，這張表的平均是 $(10 + 3 + 3 + 3 + 6)/13 = 25/13 \approx 1.92$；用期望值公式
$1 + \alpha/2 - \alpha/(2n) = 1 + 1.3 - 0.1 = 2.2$（資結書定理~12.2 第 3 點）。 兩個都不在選項裡。
\end{itemize}
所以題目問的是失敗的搜尋，答案是 $\alpha = 2.6$，選 (B)。 (D) 3.6 是 $1 + \alpha$，那是 CLRS 寫搜尋
\textbf{時間} $\Theta(1 + \alpha)$ 時多算的「計算雜湊值」那一步，不是鍵的比較。 (A) 1.0、(C) 3.0 沒有
對應的算法。

\textbf{12.} $m = 7$，依序插入：
\begin{center}
\small
\begin{tabular}{rcccl}
\toprule
$k$ & $h_1(k) = k \bmod 7$ & $h_2(k) = 5 - (k \bmod 5)$ & 探測 & 放在 \\
\midrule
35 & 0 & 5 & 0 & 0 \\
42 & 0 & 3 & 0（有 35）, $0 + 3 = 3$ & 3 \\
3 & 3 & 2 & 3（有 42）, $3 + 2 = 5$ & 5 \\
21 & 0 & 4 & 0（有 35）, $0 + 4 = 4$ & 4 \\
\bottomrule
\end{tabular}
\end{center}
$21$ 放在第 $4$ 格，選 (C)。 (A) 0 是 $h_1(21)$，被 $35$ 佔了；(B) 3 和 (D) 5 是 $42$、$3$ 的位置。
\end{pastsoln}
